% ============================================================================
% BRIEF METHODS AND RESULTS SECTION FOR MAIN PAPER (v2)
% ============================================================================
% Restructured: Baselines → Performance Comparison → Validation Method → Final Framework
% Full details in Supplementary Materials.
% ============================================================================

\documentclass[11pt,a4paper]{article}
\usepackage[utf8]{inputenc}
\usepackage{amsmath}
\usepackage{booktabs}
\usepackage{geometry}
\usepackage{cite}
\usepackage{array}
\geometry{margin=1in}

\begin{document}

\section*{Methods}

\subsection*{Baseline Methods and Performance Comparison}

We evaluated four classification approaches on 759,774 cardiovascular DDI pairs from DrugBank v5.1.9:

\textbf{(1) Zero-Shot BART-MNLI.} DDI descriptions classified via \texttt{facebook/bart-large-mnli}, a 400M-parameter transformer framing severity as textual entailment against four hypotheses. Baseline showed severe over-classification: 56.9\% Contraindicated, 43.0\% Major.

\textbf{(2) Confidence-Weighted.} Extended zero-shot with penalization for low-confidence ($<$65\%) high-severity predictions.

\textbf{(3) Rule-Based.} Clinical keyword pattern matching against severity-specific markers (e.g., ``contraindicated,'' ``life-threatening'' $\rightarrow$ Contraindicated).

\textbf{(4) Evidence-Based Classifier.} Hybrid framework incorporating FDA Black Box Warnings, ACC/AHA clinical guidelines \cite{joglar2023,january2019}, and CHEST recommendations with drug class risk profiling.

\begin{table}[h]
\centering
\small
\caption{Method Performance Comparison (DDInter Validation, n=37,164)}
\label{tab:methods}
\begin{tabular}{@{}lccc@{}}
\toprule
\textbf{Method} & \textbf{Exact Acc.} & \textbf{Adjacent Acc.} & \textbf{Cohen's $\kappa$} \\
\midrule
Zero-Shot BART-MNLI & 22.1\% & 95.6\% & 0.000 \\
Confidence-Weighted & 22.1\% & 95.6\% & 0.000 \\
\textbf{Rule-Based} & \textbf{71.2\%} & \textbf{99.9\%} & \textbf{+0.144} \\
Evidence-Based & 32.2\% & 92.3\% & +0.066 \\
\bottomrule
\end{tabular}
\end{table}

\subsection*{External Validation Method}

External validation was performed against DDInter \cite{ddinter2022}, a literature-curated DDI database with expert-validated severity annotations (Xiong et al., \textit{Nucleic Acids Research}, 2022). DDInter provides independent ground truth derived from scientific literature, package inserts, and regulatory labels---distinct from our classification methodology.

Drug name matching between DrugBank and DDInter yielded 37,164 validated pairs after filtering Unknown labels. DDInter uses three severity levels (Major/Moderate/Minor); we mapped these to our four-level scale by combining our Contraindicated and Major predictions into DDInter's Major category. Performance metrics included exact accuracy, adjacent accuracy ($\pm$1 severity level), and Cohen's $\kappa$ for inter-rater agreement.

\subsection*{Final Evidence-Based Framework}

The Evidence-Based Classifier combines three weighted components:

\begin{equation}
S_{\text{final}} = 0.45 \cdot S_{\text{semantic}} + 0.25 \cdot S_{\text{confidence}} + 0.30 \cdot S_{\text{drug\_class}}
\end{equation}

\textbf{(1) Semantic Similarity (0.45).} DDI descriptions encoded via \texttt{all-MiniLM-L6-v2} sentence embeddings compared to severity-specific prototypes representing canonical clinical scenarios.

\textbf{(2) Confidence Adjustment (0.25).} Penalizes low-confidence ($<$65\%) high-severity predictions from the base classifier.

\textbf{(3) Drug Class Risk (0.30).} Elevates severity for high-risk classes: anticoagulants, QT-prolonging agents, MAOIs, based on FDA Black Box Warnings and ACC/AHA guidelines \cite{joglar2023,beavers2022}.

A curated set of 34 DDI pairs with documented FDA Black Box Warnings bypasses scoring and receives direct severity assignment. Full methodology in Supplementary S1--S4.

\section*{Results}

The Rule-Based classifier achieved 71.2\% exact accuracy and 99.9\% adjacent accuracy against DDInter external validation (n=37,164), with Cohen's $\kappa$=+0.144---the highest inter-rater agreement among all methods (Table~\ref{tab:methods}). Both Rule-Based and Evidence-Based methods substantially outperformed zero-shot classification. Recalibration corrected baseline over-classification: Jensen-Shannon divergence decreased from 0.847 to 0.000, with 714,290 predictions (94\%) modified to match target distribution (Table~\ref{tab:dist}). GPU-accelerated processing achieved 15,454 interactions/second (759,774 DDIs in 49s). Validation details in Supplementary S7--S9.

\begin{table}[h]
\centering
\small
\caption{Severity Distribution: Original vs. Recalibrated}
\label{tab:dist}
\begin{tabular}{@{}lrrr@{}}
\toprule
\textbf{Severity} & \textbf{Original} & \textbf{Recalibrated} & \textbf{Target} \\
\midrule
Contraindicated & 56.9\% & 5.0\% & 5\% \\
Major & 43.0\% & 25.0\% & 25\% \\
Moderate & $<$0.1\% & 60.0\% & 60\% \\
Minor & 0.1\% & 10.0\% & 10\% \\
\bottomrule
\end{tabular}
\end{table}

\begin{thebibliography}{99}

\bibitem{joglar2023}
Joglar JA, Chung MK, Armbruster AL, et al.
2023 ACC/AHA/ACCP/HRS Guideline for the Diagnosis and Management of Atrial Fibrillation.
\textit{J Am Coll Cardiol}. 2024;83(1):109-279. PMID: 38043043.

\bibitem{january2019}
January CT, Wann LS, Calkins H, et al.
2019 AHA/ACC/HRS Focused Update of the 2014 AHA/ACC/HRS Guideline for the Management of Patients With Atrial Fibrillation.
\textit{J Am Coll Cardiol}. 2019;74(1):104-132. PMID: 30703431.

\bibitem{beavers2022}
Beavers CJ, Rodgers JE, Bagnola AJ, et al.
Cardio-Oncology Drug Interactions: A Scientific Statement From the American Heart Association.
\textit{Circulation}. 2022;145(15):e811-e838. PMID: 35249373.

\bibitem{ddinter2022}
Xiong G, Yang Z, Yi J, et al.
DDInter: an online drug-drug interaction database towards improving clinical decision-making and patient safety.
\textit{Nucleic Acids Res}. 2022;50(D1):D1200-D1207. PMID: 34634800.

\end{thebibliography}

\end{document}
